\section{Research Problem and Objectives}

\subsection{Research Problem Formulation}
Existing literature separately evaluates conformal uncertainty calibration, operating-condition diagnostics, and uncertainty-aware maintenance scheduling. What remains insufficiently established is a Remaining Useful Life (RUL) prediction methodology that jointly adapts uncertainty calibration to operating conditions and downstream sequential replacement consequences while preserving reliable interval coverage.

Our study investigates whether decision-aware, condition-adaptive conformal calibration can improve sequential maintenance utility without sacrificing subgroup reliability or merely becoming more conservative.

\subsection{Formal Research Questions and Hypotheses}
In accordance with the frozen research protocol (ADR-018), four prespecified research questions are evaluated under version control prior to held-out evaluation:
\begin{itemize}
  \item \textbf{RQ1 (Sequential Maintenance Utility \& Coverage)}: \textit{Can a decision-aware, condition-adaptive conformal RUL calibration method improve sequential maintenance utility while preserving reliable prediction-interval coverage?}
  Evaluated on FD002 at 90\% nominal coverage under \texttt{primary\_reference} cost: requires cost improvement $\Deltacost > 0.20~\SCU$ per engine over both global and condition-aware calibration, marginal, worst-regime, and worst-stage absolute coverage error ceilings $< 0.05$, and observed fleet failure rate increase $\Deltafail \le 0.0$.
  \item \textbf{RQ2 (Conditional Subgroup Reliability)}: \textit{Can the proposed method improve reliability across operating regimes and degradation conditions without excessively widening prediction intervals?}
  Evaluated across operating regimes ($C_1$--$C_6$) and degradation stages (Early/Mid/Late): requires worst-group conditional coverage error improvement $\Deltaerror > 0.01$ over both baselines, absolute worst-group error $< 0.05$, and width and Winkler score deteriorations $\Dwidth, \Dscore < 1.0$~cycle.
  \item \textbf{RQ3 (Decision-Value Attribution)}: \textit{Are the maintenance benefits of the proposed method attributable to decision-aware uncertainty calibration rather than simply more conservative prediction intervals or earlier maintenance?}
  Evaluated against four prespecified matched-conservatism controls across three families (matched width, widened interval, point safety margin, earlier intervention): requires cost improvement $\Deltacost > 0.20~\SCU$ over every control, proposed-method coverage error ceilings $< 0.05$, and observed failure rate increase $\Deltafail \le 0.0$.
  \item \textbf{RQ4 (Robustness \& Generalization)}: \textit{Does the proposed method retain its calibration and maintenance-decision benefits across operating conditions, base RUL predictors, datasets, and repeated experimental runs?}
  Evaluated across the balanced 6-cell matrix (FD001, FD002, FD004 $\times$ GBR, LSTM) weighted equally ($1/6$ each): requires weighted aggregate $\Deltacost > 0.20~\SCU$ and $\Deltaerror > 0.01$, cell-level non-harm ($\Deltacost > -0.20~\SCU$, $\Deltaerror > -0.01$), and all cell-level coverage, width, score, and observed-failure safeguards.
\end{itemize}
